\documentclass[11pt,a4paper]{article}
\usepackage[T1]{fontenc}
\usepackage{lmodern,amsmath,amssymb,microtype,booktabs}
\usepackage[margin=25mm]{geometry}
\usepackage[colorlinks=true,urlcolor=blue,linkcolor=blue]{hyperref}
\title{Transverse phase of the complete dS$_3$ crossing kernel:\\
a punctured-window expansion and logarithmic scaling}
\author{Alan Day \and Chengheng\thanks{AI research assistant (OpenAI).}}
\date{Working research note --- 31 October 2026}
\begin{document}\maketitle
\begin{abstract}
The physical normalization phase changes the logarithmic transverse
scaling inferred from one Omega factor alone. Using the established
Wilson reflection and the two-series LPRS formula, we obtain a
phase-shifted sine kernel away from the collision surface and a
pointwise sinc limit for the complete scalar matrix element.
An explicit finite-mass bound retains the whole companion series.
Neither result establishes a distributional limit or the heavy-mass
limit of a physical wave packet.
\end{abstract}

\section{Model, coordinates and claims}
The radius is one and $G_N=0$. We keep proper $SO^+(3,1)$, BD scalar
observers, KG matter, a fixed positive observer-label lower bound,
and half-period regulated creation blocks. $U_{00}$ is the scalar
matrix element in the physical CG-period convention, not a separate
unitary scalar transform. The principal-series label $M$ corresponds
to static energy $\sqrt{1+M^2}$.

Fix real $u,v$, positive $p,q$, and put
\[
 c=M,\quad b=M+u,\quad a=M+u+v,\quad d=M+v+r,\qquad z=-ir/2,
\]
\[
 A_j=\tfrac12+i\alpha_j,\qquad
 (\alpha_1,\alpha_2,\alpha_3,\alpha_4)
 =\tfrac12(v+q,p-u,-u-p,v-q).
\]
Our real transfer window requires
$|\alpha_j|\le1/2$ and $|\alpha_j+r/2|\le1/2$.
Define a continuous phase, with $\delta(0)=0$, by
\begin{equation}
 \delta(r)=\Im\left[
 2\log\Gamma(1-ir/2)+
 \sum_{j=1}^4\{\log\Gamma(A_j+ir/2)-\log\Gamma(A_j)\}\right].
 \label{delta}
\end{equation}
For fixed $r\ne0$, uniformly on compact subsets of this punctured window,
\begin{equation}
 U_{00}=\frac{\sin(r\log M-\delta(r))}{\pi r}+O(M^{-1}).
 \label{sine}
\end{equation}
For transfers strictly inside the $r=0$ window, each fixed real
$\lambda$ instead gives
\begin{equation}
 \lim_{M\to\infty}\frac{\pi U_{00}(r=\lambda/\log M)}{\log M}
 =\frac{\sin\lambda}{\lambda},
 \qquad \left.\frac{\sin\lambda}{\lambda}\right|_{\lambda=0}=1.
 \label{sinc}
\end{equation}
This is a pointwise scaling statement. The normalization phase must
be included: the single collision factor had the different profile
$\sin(\lambda/2)/(\lambda/2)$.

\newpage
\section{The two oscillatory branches and the physical phase}
The established reflection identity yields
$U_{00}=2\Re(C_+XY)$, where $X=\Omega(h,h')$ and
$Y=\Omega(1-h,h')$, $h=(1+ib)/2$, $h'=(1+id)/2$.
The two-series formula \cite{lprs}, with the Wilson reflection
\cite{wilson}, has the useful exact structure
\[
 X=c_AF+c_BG,\qquad Y=y_AF+y_BH .
\]
Here $F$ is the same series in both expressions:
\[
 F={}_4F_3(A-z;\,1+id,\,1-ib-z,\,1-z;\,1),
\]
\[
 H={}_4F_3(A;\,1+i(M+v)-z,\,1-ib,\,1+z;\,1),
\]
where $A-z$ denotes all four numerator parameters.
The large-parameter series is
\[
 G={}_4F_3(A+ib;\,1+i(b+M+v+r/2),\,1+ib,\,1+i(b-r/2);\,1).
\]
The coefficients of $Y$ are
\begin{align*}
 y_A&=\Gamma(z)\frac{\Gamma(1-ib)}{\Gamma(1-ib-z)}
             \frac{\Gamma(A_3-z)\Gamma(A_4-z)}
                  {\Gamma(A_3)\Gamma(A_4)},\\
 y_B&=\Gamma(-z)\frac{\Gamma(1+id)}{\Gamma(1+i(M+v)-z)}
             \frac{\Gamma(A_1)\Gamma(A_2)}
                  {\Gamma(A_1-z)\Gamma(A_2-z)}.
\end{align*}
Let $D_0=C_+c_A$, $E_0=C_+c_B$ and define
\[
 B_s=\frac{\Gamma(A_1)\Gamma(A_2)}
            {\Gamma(A_1-z)\Gamma(A_2-z)},\quad
 A_s=\frac{\Gamma(A_3-z)\Gamma(A_4-z)}
            {\Gamma(A_3)\Gamma(A_4)}.
\]
Substitution of the physical three-sphere periods gives
\begin{align}
 D_0&=D_\infty M^z[1+O(M^{-1})],\nonumber\\
 |D_\infty|&=\frac{e^{\pi r/4}}{4\pi|\Gamma(1-z)|}
 \sqrt{\prod_{j=3,4}
 \frac{\cosh\pi(\alpha_j+r/2)}{\cosh\pi\alpha_j}},\nonumber\\
 \arg D_\infty&=-\arg B_s+\arg\Gamma(1+z).
 \label{D}
\end{align}
The phases here are continuous, not principal-value fits. The Gamma
argument/weight certificate is saved in the audit. Shifted Stirling
expansion \cite{dlmf} gives the large-argument phase explicitly as
$-r\log M/2-r(3r+4v)/(8M)+O(M^{-2})$; the remaining terms give
the phase in \eqref{D}. The modulus follows from
$|\Gamma(1/2+ix)|^2=\pi/\cosh\pi x$ and
$|\Gamma(1+ix)|^2=\pi x/\sinh\pi x$.

Gamma ratios imply $y_A\sim\Gamma(z)A_s(-iM)^z$ and
$y_B\sim\Gamma(-z)B_s(iM)^{-z}$. Thus
\[
 D_0(y_A+y_B)=T(r)e^{-ir\log M}+J(r)+O(M^{-1}),
\]
\[
 T=D_\infty\Gamma(z)e^{-i\pi z/2}A_s
   =\frac{i\,e^{i\delta(r)}}{2\pi r},\qquad
 J=D_\infty\Gamma(-z)e^{-i\pi z/2}B_s\in i\mathbb R .
\]
The nonoscillating branch has zero real part at leading order.
Its cancellation is algebraic: $\Gamma(-z)\Gamma(1+z)
=z|\Gamma(z)|^2$ is imaginary. The cosh factors cancel in $|T|$,
leaving $1/(2\pi|r|)$. Taking twice the real part proves the
leading expression in \eqref{sine}, once the series errors are controlled.

\newpage
\section{A finite-mass bound and the limits it supports}
Put $S=y_A+y_B$, $m_A=\min(d,b-r/2)>0$,
$m_B=\min(M+v+r/2,b)>0$. The small-numerator coefficient
comparison used in the collision analysis gives
\[
 |F-1|\le R_F=\frac{\zeta(2)-1}{1+m_A^2},\qquad
 |Y-S|\le R_Y=(\zeta(2)-1)
 \left(\frac{|y_A|}{1+m_A^2}+\frac{|y_B|}{1+m_B^2}\right).
\]
These are exact-arithmetic analytic inequalities; the finite numerical
evaluation of their right-hand sides is not interval certified.

For completeness the whole $G$ series, not merely its first term,
is bounded as follows. Let
\[
 B=b-|r|/2>0,\qquad A_*=\max_j|\alpha_j|+|r|/2,
 \qquad b+M+v+r/2\ge B .
\]
For $k=n+1$ its term ratio obeys
\[
 |g_{n+1}/g_n|\le
 \sqrt{1+B^2/k^2}
 \left(1+\frac{2BA_*+A_*^2}{k^2+B^2}\right)^2.
\]
The sinh product \cite{product} bounds every term by
\[
 B_b=\sqrt{\frac{\sinh\pi B}{\pi B}}\,
       e^{2\pi A_*+\pi A_*^2/B}.
\]
The exact squared-ratio comparison gives a monic polynomial
$P(x)=x^{10}+\sum_{j<10}p_jx^j$. Set
$N_*=2+\lceil\max_{j<10}|p_j|\rceil$.
All translated coefficients at $N_*$ are positive by the monic
polynomial root bound, also applied to its normalized derivatives.
Consequently the term ratio is at most
$((n+1)/(n+2))^{3/2}$ for $n\ge N_*$, and
\[
 |G|\le B_G=(3N_*+3)B_b .
\]
Here $N_*=O(M^8)$ is an analytic threshold, not a numerical loop length.
The Gamma prefactors give $|E_0|=O(M^{-1/2}e^{-\pi M})$ on
compact transfers, so $|E_0|B_G=O(M^8e^{-\pi M/2})$ suffices.

Combining these estimates yields the explicit bound
\begin{equation}
 \left|U_{00}-2\Re(D_0S)\right|
 \le 2|D_0|\{R_Y+R_F(|S|+R_Y)\}
       +2|E_0|B_G(|S|+R_Y).
 \label{finite}
\end{equation}
At fixed nonzero $r$, this error is $O(M^{-2})$ plus an exponentially
small term. The Gamma coefficient expansion then proves \eqref{sine}.

For fixed $\lambda\ne0$, $r=\lambda/\log M$ gives
$|y_A|+|y_B|=O(\log M)$, while the coefficient-ratio expansions
remain uniform for the small shifts. Their errors are
$O(\log M/M)$, and \eqref{finite} contributes
$O(\log M/M^2)$ plus an exponential term.
Also $\delta(r)=O(r)$; division by $\log M$ proves \eqref{sinc}.
The $\lambda=0$ case uses the preceding peak theorem separately.

In fact
\[
 \delta'(0)=-\tfrac12\Re\{2\psi(1)-\sum_j\psi(A_j)\}.
\]
The removable $r=0$ value of the leading sine expression therefore
agrees with the logarithm and constant of the peak theorem.
This matching is not a proof of uniform error across $r=0$:
the separate-branch bound contains $1/|r|$.

\newpage
\section{Numerical checks, preserved failures and next task}
We use $(u,v,p,q)=(0.1,-0.1,0.4,0.4)$,
$M=16,64,256,1024$, and both signs of either $r=0.08$ or
$r=1/\log M$. Sixteen complete reflected-kernel evaluations use
90 decimal digits and at least 128 terms per series.
All obey \eqref{finite}, allowing separately for the saved truncation tail.
For positive $r=0.08$:
\begin{center}
\begin{tabular}{rrrr}
\toprule
$M$ & error vs.\ $2\Re(D_0S)$ & bound \eqref{finite}
 & error vs.\ sine form\\
\midrule
16 & $8.41\,10^{-4}$ & $9.22\,10^{-1}$ & $5.58\,10^{-3}$\\
64 & $6.74\,10^{-5}$ & $1.59\,10^{-3}$ & $1.47\,10^{-3}$\\
256 & $5.06\,10^{-6}$ & $1.03\,10^{-4}$ & $4.09\,10^{-4}$\\
1024 & $3.64\,10^{-7}$ & $6.69\,10^{-6}$ & $1.14\,10^{-4}$\\
\bottomrule
\end{tabular}
\end{center}
The explicit finite-mass bound controls a Gamma-exact approximation,
not the final sine expression. These errors must not be conflated.

The sinc convergence is slow and not monotone in this range.
For $\lambda=1$, $\pi U_{00}/\log M$ is approximately
$0.9225,0.9975,0.9977,0.9853$; the limiting value is
$\sin1\simeq0.84147$. The phase correction explains why peak-shape
plots at moderate mass would be misleading.

An independent Decimal/Fraction replay and symbolic Gamma-phase
check passes 203 checks. Four selected points were recomputed at
120 digits and doubled cutoff. At $M=16$ the relative changes are
about $1.63\,10^{-23}$ and $9.29\,10^{-24}$; at the two $M=1024$
scaled points they are about $1.52\,10^{-86}$ and $1.42\,10^{-86}$.
These checks do not establish interval rounding bounds or accuracy
at unsampled parameters.

The earlier collision parameter-derivative job has now exited normally.
Its twelve 100-digit comparisons and subsequent saved-arithmetic audit
pass. The earlier prototype coefficient bug, failed large-mass
acceleration, shadow amplification, and insufficient-pi audit are all
retained. Their old PDFs describe their historical checkpoints;
new completion records give the current verification status.

The next analytic task is a bound uniform through the transverse
collision, followed by oscillatory-tail control against the actual
creation amplitudes and spectral measures. Neither \eqref{sine} nor
\eqref{sinc} alone permits a delta replacement in a Bose four-point
function. Full Fock domains, finite-$G_N$ dynamics and a microscopic
de Sitter holographic dictionary remain open. The results here use
known special-function formulas in a specified model; no priority
claim is made.

\begin{thebibliography}{9}
\bibitem{lprs} J. Liu et al., \emph{$d$-dimensional SYK, AdS Loops, and
$6j$ Symbols}, \href{https://arxiv.org/html/1808.00612v3}{arXiv:1808.00612v3},
Eqs.\ (3.35)--(3.38).
\bibitem{wilson} W. Groenevelt, \emph{The Wilson function transform},
\href{https://arxiv.org/html/math/0306424v1}{arXiv:math/0306424v1},
Proposition 4.4 and Remark 4.5. No Wilson-measure unitarity is assumed.
\bibitem{dlmf} NIST DLMF,
\href{https://dlmf.nist.gov/5.11}{Section 5.11}, shifted Stirling and
Gamma-ratio asymptotics with sector restrictions.
\bibitem{product} NIST DLMF,
\href{https://dlmf.nist.gov/4.36}{Section 4.36.1}, sinh infinite product.
Sources checked 31 October 2026.
\end{thebibliography}
\end{document}